\subsection{极大元与极小元}

定义 3. 设 E 是一个有序集。元素 \(a \in E\) 称为 E 的极小元（相应地，极大元），如果关系 \(x \leqslant a\) （相应地 \(x \geqslant a\) ）蕴含 \(x = a\) .

E 的每个极小元都是关于相反序的极大元，反之亦然。

示例

\begin{enumerate}
\def\labelenumi{(\arabic{enumi})}
\item
  设 A 是一个集合。在 \(\mathfrak{P}(A)\) （按包含关系排序）中由 A 的非空子集组成的子集里，极小元就是由单个元素组成的子集。
\item
  在集合 \(\Phi(\mathbf{E},\mathbf{F})\) （由 \(\mathbf{E}\) 的子集到 \(\mathbf{F}\) （F 非空）的映射所组成的集合）中，按关系“ \(v\) 是 \(u\) 的延拓”（ \(u\) 与 \(v\) 之间）排序，极大元就是从整个 \(\mathbf{E}\) 到 \(\mathbf{F}\) 的映射。
\end{enumerate}

* (3) 在自然整数集 \textgreater{} 1 中，按 m 与 n 之间的关系“m 整除 n”排序，极小元就是素数。 *

* (4) 实数集没有极大元也没有极小元。 *

